\subsection{Симуляция по трассе}

Рассмотрим подробнее симуляцию по трассе, ранее упомянутую в разделах \ref{sec:execution_traces}, \ref{sec:dfg_collection}, \ref{sec:patched_trace_sim} и \ref{sec:trace_write}. В процессе симуляции данные из трассы пробрасываются в состояние симулятора; таким образом воспроизводится записанное в трассе исполнение и собираются необходимые данные, отсутствующие в трассе (например, значения регистров в определенных точках). Также обычно проверяется консистентность потоков управления и данных.

Для ряда рассмотренных задач применяется симуляция по модифицированной трассе. Заметим, что модифицированная трасса при этом еще не записана в выходной файл, а конструируется на лету из записей исходной трассы и подготовленных модификаций. Функционал обычной симуляции по трассе получается повторно использовать с помощью наследования и переопределения виртуальных функций. Переопределенные методы-обработчики подменяют исходные записи трассы (и/или добавляют новые), а после вызывают методы базового класса на полученных данных. Пример переопределенного метода для обработки модифицированной трассы приведен в листинге \ref{lst:sim_override}.

\begin{listing}[htbp]
    \inputminted{c++}{listings/sim_override.cpp}
    \caption{Переопределение метода-обработчика записей доступа в память}
    \label{lst:sim_override}
\end{listing}

Также класс симулятора по трассе поддерживает добавление обратных вызовов (callback), что позволяет повторно использовать его для различных симуляционных задач решения. На рис. \ref{fig:sim_uml} представлена UML-диаграмма классов симуляционного модуля.

\begin{figure}[h!]
    \centering
    \includegraphics[width=\linewidth]{pics/sim_uml.drawio.png}
    \caption{UML-диаграмма классов симуляционного модуля.}
    \label{fig:sim_uml}
\end{figure}

\FloatBarrier
